\documentclass[11pt]{amsart}
\usepackage[margin=0.9in]{geometry}
\usepackage{amsmath,amssymb,amsthm,mathtools}
\linespread{0.97}
\setlength{\parskip}{1pt}

\newtheorem{theorem}{Theorem}[section]
\newtheorem{proposition}[theorem]{Proposition}
\newtheorem{lemma}[theorem]{Lemma}
\newtheorem{corollary}[theorem]{Corollary}
\theoremstyle{definition}
\newtheorem{definition}[theorem]{Definition}
\newtheorem{remark}[theorem]{Remark}
\theoremstyle{remark}
\newtheorem*{checkpoint}{\textbf{Checkpoint}}

\newcommand{\Om}{\Omega}
\newcommand{\Ot}{\widetilde\Omega}
\newcommand{\Et}{E_{\hat t}}
\newcommand{\St}{\Sigma_{\hat t}}
\newcommand{\dT}{\delta_T}
\newcommand{\Asym}{\mathcal A}
\newcommand{\Hess}{D^2u}
\newcommand{\TF}{\bigl(\Hess+\tfrac1n I\bigr)}
\newcommand{\R}{\mathcal R}
\newcommand{\Vbar}{\mathcal V}
\newcommand{\Rn}{\mathbb R^n}
\newcommand{\Hn}{\mathcal H^{n-1}}
\newcommand{\norm}[1]{\left\lVert#1\right\rVert}
\newcommand{\abs}[1]{\left\lvert#1\right\rvert}
\DeclareMathOperator{\dist}{dist}
\DeclareMathOperator{\diam}{diam}
\DeclareMathOperator{\tr}{tr}
\newcommand{\fint}{\mathop{\mathchoice%
  {\,\rlap{$-$}\!\int}{\,\rlap{$-$}\!\int}%
  {\,\rlap{$\scriptstyle-$}\!\int}{\,\rlap{$\scriptscriptstyle-$}\!\int}}\nolimits}

\title{Two strategies for sharp quantitative Saint--Venant stability \\ via a single torsion level set}
\author{Working note (Plan 3)}
\date{\today}

\begin{document}
\begin{abstract}
We record, in full and honestly, the two strategies developed for the sharp
quantitative Saint--Venant (equivalently, torsional Faber--Krahn) inequality
$\Asym(\Om)^2\le C(n)\,\dT(\Om)$ through the device of a single torsion level
set. Strategy~A (``graph route'') is the by-hand Plan~1 idea executed with
Serrin stability: pass to a near-Serrin level set, show it is a small graph,
close perturbatively, transfer back. Strategy~B (``integrated
route'') uses Serrin \emph{directly}: Weinberger's $P$-function rigidity made
quantitative, controlling the trace-free Hessian and hence the asymmetry, with
no graph and no boundary regularity. We give the common skeleton with proofs,
both attacks with detailed sketches, the tentacle analysis (including the
precise origin of the $w^2$ escape floor that obstructs Strategy~A for
$n\ge3$), the sharpness/norm-gap discussion, the comparison with the BDV
selection principle, and an explicit checklist of every step requiring
independent verification.
\end{abstract}
\maketitle

\section{Setup}

Let $\Om\subset\Rn$, $n\ge2$, be open with finite measure; after scaling
$|\Om|=\omega_n=|B_1|$. Its \emph{torsion function} $u=u_\Om$ is the unique
minimiser of $J(v)=\tfrac12\int|\nabla v|^2-\int v$ over $H^1_0(\Om)$,
equivalently the variational solution of
\begin{equation}\label{eq:torsion}
  -\Delta u=1\ \text{ in }\Om,\qquad u=0\ \text{ on }\partial\Om,\qquad u>0\ \text{ in }\Om.
\end{equation}
The \emph{torsional rigidity} is $T(\Om)=\int_\Om|\nabla u|^2=\int_\Om u=-2\min J$.
For the ball $B_r$, $u_{B_r}(x)=\tfrac{r^2-|x|^2}{2n}$ and
$T(B_r)=\dfrac{\omega_n r^{\,n+2}}{n(n+2)}$. The Saint--Venant inequality
states $T(\Om)\le T(B)$ for the volume-matched ball $B$, with equality iff
$\Om$ is a ball. Define the scale-invariant deficit and the Fraenkel asymmetry
\[
  \dT(\Om):=\frac{T(B)-T(\Om)}{T(B)}\ge0,\qquad
  \Asym(\Om):=\min_{x_0}\frac{|\Om\triangle B(x_0)|}{|\Om|}.
\]
The target is the sharp quantitative inequality
\begin{equation}\label{eq:target}
  \boxed{\ \Asym(\Om)^2\le C(n)\,\dT(\Om)\ }\qquad(\dT(\Om)\le\dT^0(n)).
\end{equation}

\subsection*{Sharpness of the exponent}
For the ellipsoid $\mathcal E_\varepsilon=\{\sum x_i^2/a_i^2<1\}$ with
$a_1=1+\varepsilon$, $a_i=1-\varepsilon/(n-1)$ (volume fixed to first order),
one computes $\Asym(\mathcal E_\varepsilon)\simeq_n\varepsilon$ while
$\dT(\mathcal E_\varepsilon)\simeq_n\varepsilon^2$ (the first-order shape
variation is volume/translation-neutral, so the deficit is purely quadratic).
Hence $\Asym^2/\dT\to$ const: the exponent $2$ in \eqref{eq:target} is
optimal, and any proof must produce exactly the quadratic relation, not a
fractional power. This is why every step below is tracked for the
\emph{linear}-in-$\dT$ (i.e.\ quadratic-in-$\Asym$) scaling.

\section{The common skeleton}

Neither strategy argues on the rough set $\Om$. Instead one passes to a single
superlevel set $\Et=\{u>\hat t\}$ of its own torsion function, proves
\eqref{eq:target} \emph{there}, and transfers it back. The structural reason
to use a level set --- the insight of Plan~3 --- is:

\begin{quote}\itshape
$-\Delta u=1$ has real-analytic coefficients, so $u$ is real-analytic in the
interior of $\Om$, with \emph{no hypothesis on $\partial\Om$}. By Sard's
theorem, for a.e.\ $\hat t$ the level $\St=\{u=\hat t\}$ is a compact
real-analytic hypersurface with $|\nabla u|>0$. Boundary regularity of $\Om$
is never used; only interior regularity of $u$, which is free.
\end{quote}

\subsection{The exact level-set deficit identity}
By the coarea formula, $T(\Om)=\int_\Om u=\int_0^{\|u\|_\infty}m(t)\,dt$
with $m(t)=|\{u>t\}|=|E_t|$, and
$-m'(t)=\int_{\Sigma_t}|\nabla u|^{-1}\,d\Hn$, while differentiating
$\int_{E_t}(-\Delta u)=\int_{\Sigma_t}|\nabla u|$ gives
$m(t)=\int_{\Sigma_t}|\nabla u|\,d\Hn$. Combining these with the
isoperimetric inequality $P(t)\ge n\omega_n^{1/n}m(t)^{(n-1)/n}$ and
Cauchy--Schwarz $\bigl(\int_{\Sigma_t}1\bigr)^2\le
\bigl(\int|\nabla u|^{-1}\bigr)\bigl(\int|\nabla u|\bigr)$ produces, after the
Kohler-Jobin/Talenti algebra, the \emph{exact identity}
\begin{equation}\label{eq:identity}
  \dT(\Om)\ \simeq_n\ \int_0^{\|u\|_\infty}\!\bigl(D_H(t)+D_I(t)\bigr)\,w(t)\,dt,
  \qquad w(t)=\bigl(n^2\omega_n^{2/n}m(t)^{1-2/n}\bigr)^{-1},
\end{equation}
where, on a regular level (with $P(t)=\Hn(\Sigma_t)$, $\bar f_t=m(t)/P(t)$),
\begin{align}
  D_H(t)&=\Bigl(\int_{\Sigma_t}\!\tfrac{d\Hn}{|\nabla u|}\Bigr)
          \Bigl(\int_{\Sigma_t}\!|\nabla u|\,d\Hn\Bigr)-P(t)^2\ \ge0
          \quad(\text{Cauchy--Schwarz / Serrin defect}),
  \label{eq:DH}\\
  D_I(t)&=P(t)^2-n^2\omega_n^{2/n}m(t)^{2-2/n}\ \ge0
          \quad(\text{isoperimetric defect}).
  \label{eq:DI}
\end{align}
$D_H(t)=0$ iff $|\nabla u|\equiv$ const on $\Sigma_t$; $D_I(t)=0$ iff
$\Sigma_t$ is a sphere. The decisive exact reformulation is the
\emph{weighted-variance identity}
\begin{equation}\label{eq:boxed}
  \int_{\Sigma_t}\frac{\bigl(|\nabla u|-\bar f_t\bigr)^2}{|\nabla u|}\,d\Hn
  \ =\ \frac{m(t)}{P(t)^2}\,D_H(t)\ =:\ \frac{m(t)}{P(t)^2}\,\Vbar(t)\,(\dots),
\end{equation}
obtained by expanding $\int(|\nabla u|-\bar f)^2/|\nabla u|=\int|\nabla u|^{-1}\bar f^2-2\bar f P+\int|\nabla u|$ and using $\bar f P=m$.
\emph{$D_H$ is the variance of $|\nabla u|$ on $\Sigma_t$ weighted by
$1/|\nabla u|$.} This weighting is the protagonist of the whole note.

\begin{remark}[How \eqref{eq:identity} arises]
Let $\mu(t)=-m'(t)=\int_{\Sigma_t}|\nabla u|^{-1}$ and recall
$m(t)=\int_{\Sigma_t}|\nabla u|$. The Saint--Venant functional admits the
one-dimensional reduction $T(\Om)=\int_0^\infty m(t)\,dt$, and along the
torsion flow one has the differential relation
$\bigl(\tfrac{d}{dt}\bigr)\!\int_{\Sigma_t}|\nabla u|=-P(t)^2/\!\int|\nabla u|^{-1}+\dots$
Two losses appear when comparing with the ball flow: the
\emph{Cauchy--Schwarz} loss
$P(t)^2\le\bigl(\int|\nabla u|^{-1}\bigr)\bigl(\int|\nabla u|\bigr)$, whose
defect is exactly $D_H(t)$ in \eqref{eq:DH}, and the \emph{isoperimetric}
loss $P(t)^2\ge n^2\omega_n^{2/n}m^{2-2/n}$, whose defect is $D_I(t)$ in
\eqref{eq:DI}. Integrating the resulting differential inequality for the
profile gap against the weight $w(t)$ (Talenti/Kohler-Jobin) and using
$T(B)-T(\Om)=\int_0^{|\Om|}h$ collapses to \eqref{eq:identity}. The point for
us: $\dT$ is \emph{exactly} the total of the two defects, each nonnegative,
each with a clean geometric meaning; smallness of $\dT$ forces both small
\emph{in integrated average}.
\end{remark}

\subsection{Common Lemma 1: extraction of a good near-boundary level}
\begin{lemma}[good-level extraction]\label{lem:step1}
Fix a small absolute constant $\theta\in(0,1)$. There exist $C_n$ and a level
$\hat t$ with: \textup{(i)} $D_H(\hat t)\le n^2\omega_n\,\theta$, equivalently
$\Vbar(\hat t)\le\theta$; \textup{(ii)} $|\Om\setminus\Et|\le(C_n/\theta)\dT(\Om)$;
\textup{(iii)} $\hat t$ Sard-regular, so $\St$ is real-analytic, $|\nabla u|>0$.
\end{lemma}
\begin{proof}
Drop $D_I\ge0$ in \eqref{eq:identity}: $\int(D_H w)\,dt\le C_n\dT(\Om)$. Near
the boundary $m(t)\uparrow|\Om|$, so $w(t)$ is two-sidedly bounded by
dimensional constants there --- \emph{not} because $\dT$ is small but because
$m\to|\Om|$. Change variables to the volume parameter $s=m(t)\in(0,|\Om|)$,
strictly decreasing and absolutely continuous; the layer is then
\emph{exactly linear}, $|\Om\setminus E_t|=|\Om|-s$. In $s$, the deficit
integrand $\Phi(s):=D_H w\,|dt/ds|$ satisfies $\int_0^{|\Om|}\Phi\le C_n\dT$;
by Markov the set $\{\Phi>\theta'\}$ has $s$-measure $\le C_n\dT/\theta'$.
Choosing the small-layer window $S_\lambda=[|\Om|-\lambda,|\Om|)$ with
$\lambda=(4C_n/\theta')\dT$, the bad set occupies $\le\tfrac12\lambda$ of it, so
$\{D_H\le\theta\}\cap S_\lambda$ has positive $s$-measure; pull back to
positive $t$-measure (strict monotone AC change) and intersect with the
Lebesgue-null Sard set. (The naive ``$t$-length'' phrasing is only heuristic:
for a rough domain $\int|\nabla u|^{-1}$ may be unbounded, so there is no
$t$-length lower bound; the volume variable removes this entirely.)
\end{proof}
\begin{checkpoint}[C1]
$w(t)=O(1)$ on the near-boundary range with dimensional constants; the
volume-variable Markov bookkeeping. Routine.
\end{checkpoint}

\subsection{Common Lemma 2: exact deficit transfer}
\begin{lemma}[exact deficit transfer]\label{lem:step3}
Put $\Ot:=\Et$, $v:=u-\hat t$. Then $-\Delta v=1$ in $\Ot$, $v=0$ on
$\partial\Ot$: $v$ is \emph{exactly} the torsion function of $\Ot$. With
$\lambda=|\Om\setminus\Ot|$,
$\dT(\Ot)\le(1-\lambda/\omega_n)^{-1}\dT(\Om)\le2\dT(\Om)$, i.e.\
$\dT(\Ot)=\dT(\Om)+O_n(\dT(\Om)^2)$.
\end{lemma}
\begin{proof}
Write the decreasing rearrangements $u^*_\Om(s)$, $u^*_{\Ot}(s)$ and the ball
comparison profiles $b_\Om(s)=\tfrac{1-(s/\omega_n)^{2/n}}{2n}$,
$b_{\Ot}(s)=\tfrac{\tilde\rho^2-(s/\omega_n)^{2/n}}{2n}$
($\tilde\rho=(|\Ot|/\omega_n)^{1/n}$). Talenti's theorem gives the
deficit-by-profile formula $T(B_D)-T(D)=\tfrac12\int_0^{|D|}h_D(s)\,ds+\dots$
with $h_D=b_D-u^*_D\ge0$. Because $v=u-\hat t$ is \emph{exactly} the torsion
function of $\Ot$ (an identity, not a comparison), the rearrangements satisfy
$u^*_{\Ot}(s)=u^*_\Om(s)-\hat t$ for $s\in[0,|\Ot|]$, while
$b_{\Ot}(s)-b_\Om(s)=(\tilde\rho^2-1)/(2n)$ is a \emph{constant} --- the
$(s/\omega_n)^{2/n}$ terms cancel exactly. Hence $h_{\Ot}=h_\Om+\Delta$ with
$\Delta=\hat t+(\tilde\rho^2-1)/(2n)$ a constant; the exact identity
$\hat t=u^*_\Om(|\Ot|)$ together with the Talenti bound
$u^*_\Om(|\Ot|)\le b_\Om(|\Ot|)=(1-\tilde\rho^2)/(2n)$ forces $\Delta\le0$. So
the \emph{unnormalised} deficit does not increase; the only enlargement is the
volume renormalisation factor $(1-\lambda/\omega_n)^{-1}=1+O(\lambda)$. No
``deficit monotonicity under super-levelling'' (a false statement, the V13
trap) is invoked: only the exact profile identity.
\end{proof}
\begin{checkpoint}[C2]
The profile-shift identity $u^*_{\Ot}=u^*_\Om-\hat t$ and the sign
$\Delta\le0$ (the $(s/\omega_n)^{2/n}$ cancellation is the load-bearing line).
\end{checkpoint}

\subsection{Common Lemma 3: asymmetry transfer}
\begin{lemma}[superlevel asymmetry transfer]\label{lem:step5}
For measurable $\Ot\subset\Om$, $\ \Asym(\Om)\le\Asym(\Ot)+2|\Om\setminus\Ot|/|\Om|$.
\end{lemma}
\begin{proof}
Let $B$ be an optimal ball for $\Ot$, $|B|=|\Ot|$. Take the concentric
$B'\supset B$ with $|B'|=|\Om|$; then $|B'\setminus B|=|\Om|-|\Ot|=\lambda$. By
the triangle inequality for symmetric differences,
$|\Om\triangle B'|\le|\Om\triangle\Ot|+|\Ot\triangle B|+|B\triangle B'|
\le\lambda+|\Ot\triangle B|+\lambda$. Dividing by $|\Om|$ and recalling
$\Asym(\Ot)=|\Ot\triangle B|/|\Ot|$ with $|\Ot|\le|\Om|$ gives the claim with
constant $2$.
\end{proof}
Composing Lemmas \ref{lem:step1}--\ref{lem:step5}: if at the level one has
$\Asym(\Et)^2\le c(n)\dT(\Et)$, then with $\lambda=O(\dT)$,
\[
  \Asym(\Om)^2\le2\Asym(\Et)^2+8(\lambda/|\Om|)^2
  \le2c(n)\dT(\Et)+C'(n)\dT(\Om)^2\le C_*(n)\,\dT(\Om),
\]
the squared collar $O(\dT^2)\ll\dT$ being the ``$\dT$ of room to spare''.
\begin{checkpoint}[C3]
The constant $2$ above, and the composition of volume rescalings across
Lemmas \ref{lem:step3}--\ref{lem:step5} (no $O(\dT)$ term escaping to
lower order than $\dT$).
\end{checkpoint}

\emph{Everything thus reduces to the single at-level inequality}
$\Asym(\Et)^2\le c(n)\dT(\Et)$. The two strategies are two proofs of it.

\section{Strategy A: the graph route}

This is the by-hand Plan~1 idea with Serrin stability as the regulariser:
exhibit a nearby domain that is a small \emph{graph} over a sphere and whose
asymmetry controls the original's, then invoke the sharp inequality in the
nearly-spherical class. Here the nearby domain is $\Et$ itself.

\subsection{The intended argument}
At $\hat t$ we have $D_H(\hat t)\le\theta$, i.e.\ via \eqref{eq:boxed}
$|\nabla u|$ is $L^2(\St)$-close to a constant: the \emph{approximate Serrin
condition}. Serrin's theorem: if $|\nabla v|\equiv c$ on $\partial\Ot$ for the
torsion function $v$ of $\Ot$, then $\Ot$ is a ball. Quantitative Serrin
stability (Aftalion--Busca--Reichel; Brandolini--Nitsch--Salani--Trombetti;
Magnanini--Poggesi) should upgrade ``$|\nabla v|$ near constant'' to
``$\partial\Ot$ is a small graph $\{(1+\varphi(\xi))\xi\}$ over a sphere''.
Then the sharp nearly-spherical Saint--Venant inequality (Fuglede; Brasco--De
Philippis--Velichkov) gives $\Asym(\Et)^2\lesssim\dT(\Et)$.

\subsection{What the route structurally requires}
To pass from the $L^2$ defect to the \emph{pointwise/geometric} conclusion one
needs, $\dT$-uniformly:
\begin{itemize}
\item[\textbf{(NDg)}] $|\nabla u|\ge q_->0$ \emph{everywhere} on $\St$: needed
to parametrise $\{u=\hat t\}$ as a graph at all (the implicit function theorem
divides by the normal derivative $|\nabla u|$; where $|\nabla u|\to0$ the level
can pinch or fold), and because the principal curvatures of $\St$ are
$\sim|\Hess|/|\nabla u|$.
\item[\textbf{(Mod)}] a $\dT$-uniform $C^\alpha(\St)$ modulus of $|\nabla u|$,
to run the surface interpolation
$\norm{f}_{L^\infty(\St)}\lesssim\norm{f}_{L^2(\St)}^{\vartheta}[f]_{C^\alpha(\St)}^{1-\vartheta}$,
$\vartheta=\tfrac{2\alpha}{2\alpha+n-1}$, with $f=|\nabla u|-\bar f_{\hat t}$,
$\norm{f}_{L^2}^2\simeq D_H$.
\end{itemize}
The interpolation inequality is the standard ball-covering one: at the
$L^\infty$-point $x_0$, on $\St\cap B_r(x_0)$ one has
$|f|\ge\norm f_\infty-[f]_{C^\alpha}r^\alpha$; choosing
$r=(\norm f_\infty/2[f]_{C^\alpha})^{1/\alpha}$ and using the lower volume
bound $\Hn(\St\cap B_r)\ge c r^{n-1}$ (a non-degeneracy of $\St$ itself!) gives
the displayed estimate. Both (NDg) and (Mod) thus enter unavoidably.

\subsection{The tentacle and the precise origin of the $w^2$ floor}
\label{ss:tent}
This is the point flagged as ``weird''. We derive the $w^2$ transparently.

\textbf{Model.} $\Om=B_1\cup T$, $T$ a tube of axial length $L=O(1)$ and
cross-radius $w\ll1$ attached to $\partial B_1$; $|T|\simeq\kappa_n w^{n-1}L$.

\emph{(a) Capping height $t_{\mathrm{cap}}\simeq_n w^2$.}
Freeze the tube axis; the cross-section is an $(n-1)$-ball of radius $w$.
To leading order $u$ restricted to a cross-section solves
$-\Delta'\phi=1$, $\phi|_{\partial}=0$, with the explicit barrier
$\phi(r)=\tfrac{w^2-r^2}{2(n-1)}$, $\max\phi=\tfrac{w^2}{2(n-1)}$. The
transverse Dirichlet eigenvalue is $\lambda_1^\perp\simeq w^{-2}$, so by a
Phragm\'en--Lindel\"of/maximum-principle barrier the axial direction adds only
a bounded correction: $\sup_T u\le C_n w^2$ \emph{independently of $L$}, and
$\ge c_n w^2$. Hence
\[
  t_{\mathrm{cap}}:=\sup_T u\ \simeq_n\ w^2,
\]
and for $t>t_{\mathrm{cap}}$ the level set $E_t$ does not meet $T$.

\emph{(b) The unavoidable ball-shell, measure $\simeq_n w^2$.}
A tentacle-free level needs $\hat t>t_{\mathrm{cap}}\simeq w^2$, hence one
discards $\{u<\hat t\}$. Besides $T$, this contains the boundary layer of the
ball part: with $u_{B_1}(x)=\tfrac{1-|x|^2}{2n}$,
\[
  \{x\in B_1:u_{B_1}(x)<\hat t\}=\{|x|>\sqrt{1-2n\hat t}\}
  =\{|x|>1-n\hat t+O(\hat t^2)\},
\]
an annular shell of thickness $\simeq n\hat t\simeq_n w^2$, of measure
\[
  \bigl|\{x\in B_1:u_{B_1}<\hat t\}\bigr|\simeq_n\Hn(\partial B_1)\cdot n\hat t
  \simeq_n w^2 .
\]
\emph{This is the $w^2$.} It is elementary: escaping a tube of height
$\simeq w^2$ forces peeling a ball-shell of the same height, whose
$n$-volume is $\simeq w^2$ \emph{for every dimension}, irrespective of the
(possibly tiny) tube volume.

\emph{(c) The deficit is only $\simeq_n|T|$.}
Energy expansion: attaching the thin tube changes $T(\Om)$ by
$\int_T u+(\text{junction})\simeq w^2|T|=O(w^{n+1}L)$ (negligible), while the
equal-volume comparison ball $B_*$ grows: $T(B_*)=T(B_1)(1+|T|/\omega_n)^{(n+2)/n}
=T(B_1)\bigl(1+\tfrac{n+2}{n}\tfrac{|T|}{\omega_n}+\dots\bigr)$. Hence
$T(B_*)-T(\Om)\simeq_n|T|$ and
\[
  \dT(\Om)\ \simeq_n\ |T|\ =\ \kappa_n w^{n-1}L \qquad(\text{no }w^2\text{ floor}).
\]
(Independently confirmed by a direct level-set energy computation:
$\dT=(2n^2\omega_n)^{-1}|T|(1+o(1))$.)

\emph{(d) The obstruction to Strategy~A.}
The escape layer is $\lambda_{\mathrm{esc}}\simeq_n|T|+w^2$, while the route's
budget (Lemma~\ref{lem:step1}\,(ii)) is $|\Om\setminus\Et|=O(\dT)=O(|T|)$. For
$n\ge3$, $L=O(1)$, $w\to0$ the shell dominates, $w^2\gg w^{n-1}L=|T|\simeq\dT$,
so $\lambda_{\mathrm{esc}}/\dT\simeq w^{-(n-3)}\to\infty$: a tentacle-free
(graph-like) level is \emph{unreachable within an $O(\dT)$ layer for $n\ge3$}.
Only $n=2$ survives ($|T|=\kappa_1 wL\gg w^2$ since $L\gg w$). This is exactly
``tentacles yield non-graphs''.

\subsection{Why this is not fatal to the program}
The ``weird'' feeling is justified: the obstruction is to the \emph{graph
formulation only}, and the offending configuration is excluded by the very
hypothesis $D_H(\hat t)\le\theta$. A genuine tentacle has
$|\nabla u|\simeq w\to0$ on the tube part of $\Sigma_t$ for $t<t_{\mathrm{cap}}$.
But by \eqref{eq:boxed} $D_H$ is the $1/|\nabla u|$-\emph{weighted} variance:
on the tube part $|\nabla u|\simeq w$ is far from the mean $\bar f\simeq1$ set
by the dominant ball part, so the integrand is
$(|\nabla u|-\bar f)^2/|\nabla u|\simeq 1/w$. Explicitly, decompose the level
$\Sigma_t=\Gamma_{\mathrm{ball}}\sqcup\Gamma_{\mathrm{tube}}$ for
$t<t_{\mathrm{cap}}$: on $\Gamma_{\mathrm{ball}}$, $|\nabla u|\simeq_n1$ and
$\Hn(\Gamma_{\mathrm{ball}})\simeq_n1$; on $\Gamma_{\mathrm{tube}}$,
$|\nabla u|\simeq w$ and $\Hn(\Gamma_{\mathrm{tube}})\simeq_n w^{n-2}L$. The
two-block Cauchy--Schwarz deficit is exactly
\[
  D_H(t)=A_0A_1\frac{(g_0-g_1)^2}{g_0\,g_1},\quad
  A_i=\Hn(\Gamma_i),\ \ g_i=\text{mean of }|\nabla u|\text{ on }\Gamma_i,
\]
(the value of $\bigl(\int|\nabla u|^{-1}\bigr)\bigl(\int|\nabla u|\bigr)-P^2$
for a two-valued $|\nabla u|$). With $g_0\simeq1$, $g_1\simeq w$,
$A_0\simeq1$, $A_1\simeq w^{n-2}L$:
\[
  D_H(t)\ \simeq_n\ 1\cdot w^{n-2}L\cdot\frac{1}{w}\ =\ w^{\,n-3}L
\]
on a $t$-window of length $\simeq w^2$ (reproduced exactly by
Lemma~\ref{lem:B4}). The dimensional behaviour of $w^{\,n-3}L$
(thin tube, $w\to0$, $L=O(1)$) is decisive and must be read correctly:
\begin{itemize}
\item $n=2$: $D_H\simeq w^{-1}L\to\infty$ --- a level with $D_H\le\theta$
  \emph{cannot} pass through a tentacle; the tentacle is genuinely
  priced out.
\item $n=3$: $D_H\simeq L=O(1)$ --- bounded; excluded only if the
  absolute constant exceeds $\theta$, i.e.\ a constant-dependent
  borderline, not a robust exclusion.
\item $n\ge4$: $D_H\simeq w^{\,n-3}L\to0$ --- the tentacle has
  \emph{small} $D_H$ and therefore \emph{satisfies} the near-Serrin
  hypothesis $D_H\le\theta$. It is \textbf{not} excluded.
\end{itemize}
\emph{Correction of an earlier over-claim.} A previous version of this
note asserted ``$D_H\le\theta$ forbids a tentacled level for $n\ge3$,
so the obstruction is vacuous''. That is \textbf{false for $n\ge4$}:
$w^{n-3}L\to0$. For $n\ge4$ a thin tentacle is a genuine
near-Serrin level that is not a graph and cannot be escaped within an
$O(\dT)$ layer --- Strategy~A is \emph{genuinely obstructed for
$n\ge4$}, not vacuously. What survives is \emph{Strategy~B}, which
never needs to exclude the tentacle: a tentacled level still has
$\Asym\simeq|T|$ and $\dT\simeq|T|$, so $\Asym^2\simeq|T|^2\ll|T|\simeq\dT$
and the interior estimate holds with room. (The first agentic
$L^2\!\to\!L^\infty$ attempt did err by treating $D_H$ as the
\emph{unweighted} variance; but the corrected weighted reading prices
out tentacles only for $n=2$, not for $n\ge3$.)

\begin{checkpoint}[A1, the one you doubt --- revised]
The chain (a)--(d): $t_{\mathrm{cap}}\simeq w^2$ (transverse barrier
dominates, $L$ irrelevant); ball-shell measure $\simeq_n w^2$;
$\dT\simeq_n|T|$ (no $w^2$ floor in the deficit); hence the $n\ge3$
graph obstruction. The $D_H\simeq w^{n-3}L$ tube-level scaling excludes
tentacles \emph{only for $n=2$}; for $n\ge4$ tentacles satisfy
$D_H\le\theta$, so Strategy~A's obstruction is \emph{real and not
vacuous} for $n\ge4$. Strategy~A is therefore genuinely a graph
obstruction in dimension $\ge4$; only the interior route (Strategy~B)
survives it.
\end{checkpoint}

\subsection{Honest status of Strategy~A}
Even granting (NDg) (delivered, in measure, by $D_H$ via Lemma~\ref{lem:B4}),
two genuine difficulties remain. First, the available sharp nearly-spherical
\emph{torsional} inequality (Brasco--De Philippis--Velichkov) is, as
literally proved, conditional on
$\norm{\varphi}_{C^{2,\gamma}}$ small, whereas the torsional deficit controls
only $\norm{\varphi}_{H^{1/2}}$ (see \S\ref{sec:norm}); the regularity-free
$W^{1,\infty}$ torsional analogue is not in the corpus. Second, (Mod): even
with (NDg), the interpolation needs a $\dT$-uniform $C^\alpha/\Hess$ bound on
$\St$, while interior Schauder gives only a constant
$\sim\dist(\St,\partial\Om)^{-\beta}$ that degrades as the level nears a rough
$\partial\Om$.
\begin{checkpoint}[A2]
Whether the perturbative theorem runs with the weak norm the deficit controls
($H^{1/2}$, not $C^{2,\gamma}$), and whether (Mod) can be made $\dT$-uniform.
\emph{Open for Strategy~A; precisely what Strategy~B circumvents.}
\end{checkpoint}

\section{Strategy B: a direct interior approach via Weinberger rigidity}

Serrin's theorem states that if the torsion function of a domain has constant
gradient on the boundary, the domain is a ball. Weinberger's proof of this is
not geometric: it produces an auxiliary function whose failure to be harmonic
is exactly the trace-free part of $\Hess$, and that quantity vanishes
precisely on the ball quadratic. Strategy~B makes this rigidity
\emph{quantitative} and applies it, not to the rough set $\Om$, but to a
single superlevel set of $u$, which is automatically smooth. Because the whole
argument is interior, it uses \emph{no} regularity of $\partial\Om$ and never
represents a level set as a graph.

\subsection{The Weinberger function (classical)}
Let $u$ solve \eqref{eq:torsion} on a domain $D$. Set
$P:=|\nabla u|^2+\tfrac2n u$. Since $\Delta u\equiv-1$,
\[
  \Delta P=\Delta|\nabla u|^2+\tfrac2n\Delta u
  =2|\Hess|^2-\tfrac2n
  =2\Bigl(|\Hess|^2-\tfrac1n\Bigr).
\]
The pointwise Cauchy--Schwarz inequality
$|\Hess|^2\ge(\tr\Hess)^2/n=(\Delta u)^2/n=1/n$ gives $\Delta P\ge0$, with
equality at a point iff $\Hess=-\tfrac1n I$ there, i.e.\ iff $u$ agrees to
second order with a ball quadratic $a-|x-x_0|^2/(2n)$. Writing the trace-free
Hessian $\TF$, one has $|\TF|^2=|\Hess|^2-\tfrac1n\ge0$ and the exact identity
\begin{equation}\label{eq:DP}
  \Delta P=2\,\bigl|\TF\bigr|^2 .
\end{equation}
Thus $\R(D):=\int_D|\TF|^2$ measures how far $u$ is from a ball quadratic; it
vanishes iff $D$ is a ball. Serrin/Weinberger rigidity is the qualitative
statement ``$\R(D)=0\Rightarrow D$ a ball''. Strategy~B is its quantitative
form. This subsection is classical and not in question.

\subsection{Step 1: reduction to one level set}
By the common skeleton (Lemmas~\ref{lem:step1}--\ref{lem:step5}) it suffices
to prove the sharp inequality on a single superlevel set. Precisely, there is
a level $\hat t$ such that, with $\Ot:=\Et=\{u>\hat t\}$ and
$\St=\{u=\hat t\}$:
\textup{(i)} $\St$ is a compact real-analytic hypersurface with
$|\nabla u|>0$ on it (Sard's theorem and interior analyticity; no boundary
regularity);
\textup{(ii)} the boundary defect is a small absolute constant,
$\int_{\St}(|\nabla u|-\bar f)^2/|\nabla u|\,d\Hn\le\theta$, where
$\bar f$ is the mean of $|\nabla u|$ on $\St$;
\textup{(iii)} $|\Om\setminus\Ot|\le C_n\dT(\Om)/\theta$; and
\textup{(iv)} $\dT(\Ot)\le2\,\dT(\Om)$.
Moreover $u-\hat t$ is \emph{exactly} the torsion function of $\Ot$. So it
suffices to prove $\Asym(\Ot)^2\le C(n)\,\dT(\Ot)$, which the next two steps
establish.

\subsection{Step 2: the deficit dominates the Hessian defect}
\begin{proposition}\label{prop:B1}
There is a dimensional constant $c(n)>0$ such that, for every bounded domain
$D$ with the regularity of a Sard-regular level set \textup{(}in particular
$D=\Ot$\textup{)},
\[
  \dT(D)\ \ge\ c(n)\,|D|^{-1}\,\R(D).
\]
The constant depends only on $n$ --- not on any regularity, inradius, or shape
constant of $D$ --- and is additive over connected components.
\end{proposition}
\begin{proof}[Idea]
Two classical identities for $v$ the torsion function of $D$
($v=0$, $\nabla v=v_\nu\nu$ on $\partial D$). First, the Rellich--Pohozaev
identity, obtained by testing $-\Delta v=1$ with $x\!\cdot\!\nabla v$ and
integrating by parts:
\[
  \tfrac12\!\int_{\partial D}|\nabla v|^2(x\!\cdot\!\nu)\,d\Hn
  =\tfrac{n+2}{2}\!\int_D v=\tfrac{n+2}{2}\,T(D).
\]
Second, integrating \eqref{eq:DP} and using the boundary value
$P=|\nabla v|^2$ on $\partial D$ writes the Saint--Venant gap
$T(B_D)-T(D)$ \emph{both} as a nonnegative boundary functional of
$|\nabla v|$ on $\partial D$ \emph{and}, pairing \eqref{eq:DP} with the Green
function of $D$ (a positive weight requiring no boundary regularity), as a
dimensional multiple of $\int_D|\TF|^2=\R(D)$. The pointwise inequality
$|\Hess|^2\ge1/n$, sharp only on the ball quadratic, makes the two
expressions comparable. Combining them yields the stated inequality with
$c(n)$ dimensional. This is precisely the quantitative form of Weinberger's
rigidity.
\end{proof}
\begin{checkpoint}[B1]
That Proposition~\ref{prop:B1} holds with a \emph{purely dimensional} constant,
using only that the Rellich--Pohozaev and Green-function identities are valid
on a Sard-regular interior-analytic level (no $\partial\Om$ regularity), and
that it is stable under disconnection and thin tentacles. Consistency check:
on a ball-with-thin-tube model both sides scale like the tube volume.
\end{checkpoint}

\subsection{Step 3: the Hessian defect dominates the asymmetry}
The engine is a rigidity for the Hessian.
\begin{lemma}[quantitative Liouville for the Hessian]\label{lem:B2}
Let $D\subset\Rn$ be open with $|D|=\omega_n$ and $w\in H^2(D)$. There exist a
quadratic $q(x)=a-|x-x_0|^2/(2n)$ and a constant $K(D)$ with
\[
  \int_D|\nabla w-\nabla q|^2\ \le\ K(D)\int_D\bigl|D^2w+\tfrac1n I\bigr|^2 .
\]
$K(D)$ is a Poincar\'e-type constant of $D$; it is dimensional for $D$ in a
controlled class. Controlling $K(D)$ for the level domains that arise is the
load-bearing analytic point of Strategy~B.
\end{lemma}
\begin{proposition}\label{prop:B3}
With $D=\Ot$ there is a dimensional $C(n)$ with
$\ \Asym(D)^2\le C(n)\,|D|^{-1}\,\R(D)$.
\end{proposition}
\begin{proof}[Idea]
Apply Lemma~\ref{lem:B2} with $w=u$: if $|D|^{-1}\R(D)$ is small then
$\Hess\approx-\tfrac1n I$ in $L^2$, so $u$ is $H^2$-close to a quadratic $q$;
its superlevel set $\{u>\hat t\}$ is then close in measure to the ball
$\{q>\hat t\}$ at the sharp (quadratic) rate, and the Fraenkel asymmetry,
being the squared $L^2$ size of the boundary deviation, is controlled by
$|D|^{-1}\R(D)$. The one delicate point is that $K(D)$ in
Lemma~\ref{lem:B2} could degrade where $\nabla u$ is small (the level set can
pinch there). Remove the set $Z=\{|\nabla u|\le s_0\}$ at a dimensional
threshold $s_0\simeq_n1/n$ and bound its measure directly:
\begin{lemma}[the small-gradient set is small]\label{lem:B4}
On $\St$, for every $s>0$,
$\ \Hn\bigl(\{|\nabla u|\le s\}\cap\St\bigr)\le
4s\,\bar f^{-2}\int_{\St}(|\nabla u|-\bar f)^2/|\nabla u|\,d\Hn.$
\end{lemma}
\noindent(Where $|\nabla u|\le s\le\bar f/2$ one has
$(|\nabla u|-\bar f)^2/|\nabla u|\ge\bar f^2/(4s)$; integrate over that set
only --- Chebyshev applied to the defect integral of (ii).) By (ii) the right
side is a small absolute constant, so $|Z|$ is small and is paid for
\emph{additively}: removing a piece changes the asymmetry by at most its
relative measure. On the remaining set $D\setminus Z$ one has
$|\nabla u|\ge s_0$, so $K(D\setminus Z)$ is dimensional, and the divergent
contribution that would come from the pinch region has been excised rather
than multiplied against surviving terms. If $D$ is disconnected the argument
runs per component; a component of order-one asymmetry has order-one deficit
by Lemma~\ref{lem:B5}, hence lies outside the small-deficit regime, so the
per-component treatment loses nothing.
\end{proof}
\begin{lemma}[strict-stability threshold]\label{lem:B5}
There are $\eta_0(n),c_0(n)>0$ with: any open $D$, $|D|=\omega_n$,
$\Asym(D)\ge\eta_0$ satisfies $\dT(D)\ge c_0$. \textup{(}$\dT$ is lower
semicontinuous and positive off balls on the compact-modulo-translation set
$\{\Asym\ge\eta_0\}$.\textup{)}
\end{lemma}
\begin{checkpoint}[B2, decisive]
That $K(D\setminus Z)$ in Lemma~\ref{lem:B2} is \emph{dimensional} uniformly
over the admissible level domains --- possibly pinched or multi-component ---
once the small-gradient set $Z$ is excised. Concretely, on a ball-with-tube
model the true values are $\Asym\simeq|T|$ and $\R\simeq|T|$, so
Proposition~\ref{prop:B3} must give $\Asym^2\simeq|T|^2\le C\,|T|$ with $C$
dimensional and room to spare. Everything else in Step~3 is elementary or is
Lemma~\ref{lem:B4}; \emph{this is the single point on which the strategy
rests.}
\end{checkpoint}

\subsection{Step 4: conclusion}
\begin{theorem}\label{thm:B6}
\textup{(}Conditional on Checkpoints \textup{B1, B2.}\textup{)} There is
$C_*(n)$ such that every finite-measure open $\Om\subset\Rn$ with
$\dT(\Om)\le\dT^0(n)$ satisfies $\Asym(\Om)^2\le C_*(n)\,\dT(\Om)$, sharp
exponent.
\end{theorem}
\begin{proof}
On $\Ot=\Et$, Propositions \ref{prop:B3} and \ref{prop:B1} give
$\Asym(\Ot)^2\le C(n)|\Ot|^{-1}\R(\Ot)\le (C(n)/c(n))\,\dT(\Ot)$, and
$\dT(\Ot)\le2\dT(\Om)$ by Step~1(iv). Lemma~\ref{lem:step5} transfers:
$\Asym(\Om)^2\le2\Asym(\Ot)^2+8(|\Om\setminus\Ot|/|\Om|)^2
\le C_*(n)\dT(\Om)$, the squared collar being $O(\dT^2)\ll\dT$ by
Step~1(iii). No boundary regularity, no graph, and no perturbative
nearly-spherical step is used.
\end{proof}

\section{Sharpness and the norm gap}\label{sec:norm}

Why does Strategy~A stall where B does not? Expand a nearly-spherical
$\partial D=\{(1+\varphi(\xi))\xi:\xi\in\mathbb S^{n-1}\}$ in spherical
harmonics $\varphi=\sum_{k\ge0}\sum_m \hat\varphi_{k,m}\,Y_{k,m}$,
$-\Delta_{\mathbb S^{n-1}}Y_{k,m}=k(k+n-2)Y_{k,m}$. The volume constraint
$|D|=\omega_n$ pins the mode $k=0$ at second order
($\hat\varphi_0=-\tfrac1{2}\sum_{k\ge1}|\hat\varphi_k|^2/\dots$), and choosing
the optimal centre kills the modes $k=1$ (translations are
asymmetry-neutral). The Fuglede second variation of the torsional energy is
\[
  \dT(D)\ \simeq_n\ \sum_{k\ge2}\gamma_k\,|\hat\varphi_k|^2,\qquad
  \gamma_k=\frac{(k-1)(2k+n-2)+\dots}{\,k+n-2\,}\ \asymp\ k,
\]
so the surviving quadratic form is comparable to
$\sum_{k\ge2}k\,|\hat\varphi_k|^2\simeq\norm{\varphi}_{H^{1/2}(\mathbb S^{n-1})}^2$
(the Steklov/Dirichlet-to-Neumann scale), \emph{not}
$\sum k^2|\hat\varphi_k|^2\simeq\norm{\varphi}_{H^1}^2$. The gap
$\gamma_k>0$ for $k\ge2$ is the quantitative Saint--Venant rigidity; its
\emph{linear} (not quadratic) growth in $k$ is the precise reason the deficit
controls only a half-derivative of the graph. There is an explicit family
$\Om_h$ with $\norm{\varphi_h}_{H^{1/2}}$ bounded but
$\norm{\varphi_h}_{H^1}\to\infty$: \emph{the deficit cannot control more than
$H^{1/2}$ of the graph.} The BDV nearly-spherical theorem, as written, needs
$\norm{\varphi}_{C^{2,\gamma}}$ small (its Hadamard-flow remainder, not its
$H^{1/2}$ Steklov core, forces the two extra derivatives). So Strategy~A must
either supply $C^{2,\gamma}$ smallness regularity-free (the dead (BReg)
route) or reprove BDV in the flow-free $H^{1/2}$ norm. Strategy~B sidesteps
this entirely: Propositions~\ref{prop:B1} and \ref{prop:B3} never see a graph
norm; they trade through the interior $H^2$ quantity $\R$, for which the sharp
quadratic relation is intrinsic (the $|\Hess|^2\ge1/n$ Newton inequality is
exactly order $2$).

\section{The Faber--Krahn eigenvalue version}

The title says Faber--Krahn; we record how the torsional statement
\eqref{eq:target} delivers the eigenvalue one. Let $\lambda_1(\Om)$ be the
first Dirichlet eigenvalue, $\lambda_1(B)$ its ball value at equal volume, and
$\delta_\lambda(\Om)=\bigl(\lambda_1(\Om)-\lambda_1(B)\bigr)/\lambda_1(B)$ the
Faber--Krahn deficit. Two routes connect the two functionals.
\begin{itemize}
\item \emph{Kohler-Jobin.} The Kohler-Jobin inequality
$T(\Om)\,\lambda_1(\Om)^{(n+2)/2}\ge T(B)\,\lambda_1(B)^{(n+2)/2}$, sharp on
balls, gives $\dT(\Om)\le c(n)\,\delta_\lambda(\Om)$ to first order: a small
eigenvalue deficit forces a small torsional deficit, so \eqref{eq:target}
yields $\Asym(\Om)^2\le C(n)\,\delta_\lambda(\Om)$ --- the sharp quantitative
Faber--Krahn inequality --- with no extra work.
\item \emph{Direct analogue.} The first eigenfunction $u_1>0$ solves
$-\Delta u_1=\lambda_1 u_1$; the Weinberger $P$-function for the eigenvalue
problem, $P_\lambda=|\nabla u_1|^2+\lambda_1 u_1^2$, is again subharmonic with
$\Delta P_\lambda=2|D^2u_1+\tfrac{\lambda_1}{n}u_1 I|^2\ge0$ and the same
trace-free Hessian rigidity. Strategy~B transfers verbatim with $\R$ replaced
by $\int|D^2u_1+\tfrac{\lambda_1}{n}u_1I|^2$ and the level sets of $u_1$ in
place of those of $u$; the level set $\{u_1>\hat t\}$ has its own eigenfunction
$u_1-\hat t$ only after the standard substitution, so the cleaner statement
goes through Kohler-Jobin. Either way the eigenvalue version is a corollary,
not a separate problem.
\end{itemize}
\begin{checkpoint}[FK]
The Kohler-Jobin reduction constant and its sharpness on balls (classical;
see Brasco--De Philippis--Velichkov). Only needed if the eigenvalue
statement is wanted directly rather than as a corollary of \eqref{eq:target}.
\end{checkpoint}

\section{Comparison with Plan 1, and what to trust}

\subsection{Plan 1 vs.\ the two strategies}
Plan~1 (the BDV \emph{selection principle}) replaces an arbitrary
near-extremal $\Om$ by a penalised minimiser $\Ot$ that is, by free-boundary
regularity, a smooth nearly-spherical set, with deficit not worse and
asymmetry not smaller, then closes by the nearly-spherical inequality. Its hard
core is the selection step's regularity (the $C^{1,\gamma}$/free-boundary
constant tracking, Plan~1's Lemmas 6--7). \emph{Strategy~A is the literal
analogue} --- the smooth competitor is a torsion level set, Serrin stability
the regulariser --- and inherits both a regularity gap and the $n\ge3$
tentacle/$w^2$ obstruction to graphicality. \emph{Strategy~B replaces the
selection principle outright}: it never produces a smooth competitor; it
applies quantitative Weinberger directly to the level set, so no graph, no
$C^{2,\gamma}$, no selection principle --- the regularity-free
small-gradient bound (Lemma~\ref{lem:B4}, driven by $D_H$) supplies the only
non-degeneracy the argument needs.
Strategy~B is the realisation of ``do Plan~1's idea, but with Serrin,
directly''.

\subsection{Consolidated checklist}
\textbf{Common skeleton.} C1 (Lemma~\ref{lem:step1}); C2
(Lemma~\ref{lem:step3}, profile shift \& sign); C3 (Lemma~\ref{lem:step5} \&
rescaling composition).\\
\textbf{Strategy~A.} A1 (\S\ref{ss:tent}: $t_{\mathrm{cap}}\simeq w^2$;
ball-shell $\simeq_n w^2$; $\dT\simeq_n|T|$; $n\ge3$ obstruction; the
$D_H\simeq w^{n-3}L$ tube scaling excludes tentacles \emph{only for
$n=2$} --- \emph{the $w^2$ you asked about, derived in
\S\ref{ss:tent}(b)}; for $n\ge4$ tentacles satisfy $D_H\le\theta$ and the
graph obstruction is genuine). A2 (perturbative theorem in the weak
norm $H^{1/2}$; $\dT$-uniform (Mod)) --- \emph{open; B circumvents it}.\\
\textbf{Strategy~B (decisive).} B1 (Prop.~\ref{prop:B1}: regularity-free
quantitative Weinberger, dimensional constant, stable under
disconnection/tentacles). B2 (Prop.~\ref{prop:B3} via Lemma~\ref{lem:B2}:
$H^2$-rigidity $\Rightarrow$ sharp Fraenkel, dimensional Poincar\'e constant
on pinched/disconnected $D$ after excising the small-gradient set ---
\emph{the single most important line}). B3 (Lemma~\ref{lem:B4} \& its
agreement with the tentacle $D_H$ scaling).

\subsection{What is, and is not, established}
Lemmas \ref{lem:step1}--\ref{lem:step5} are routine (C1--C3 bookkeeping).
Strategy~A's $n\ge3$ obstruction is, modulo A1, correct arithmetic
about tentacled domains; the $D_H\simeq w^{n-3}L$ scaling prices
tentacles out only for $n=2$, so for $n\ge4$ the obstruction is
\emph{genuine, not vacuous} --- a thin tentacle is a bona fide
near-Serrin level that is not a graph. It remains irrelevant to
Strategy~B, which never excludes the tentacle: such a level still has
$\Asym^2\simeq|T|^2\ll|T|\simeq\dT$, so the interior estimate holds with
room. Strategy~B's analytic inputs B1, B2 have been derived and passed an
internal adversarial verification with an independent cross-check (B3
reproduces the tentacle scaling exactly), but are \textbf{not} human-refereed.
Honest status: the program plausibly closes \eqref{eq:target} for all $n$ via
Strategy~B, conditional on B1 and especially B2 surviving a full referee;
Strategy~A, as a graph route, is genuinely obstructed for $n\ge3$ and is the
heuristic motivating B, not a proof.
\end{document}
